\section{Per-suite and per-dataset tables}\label{app:tables}

The tables in this appendix were generated by \texttt{make\_tables.py} (in this paper's folder) from the study's result files, with the study's own functions, and checked against its markdown tables. Toy entries are $\nstar$ at $d=30$; ``never'' means test error 0.1 was not reached on the grid $n\le16384$. ``portfolio'' is \code{hyb_portfolio} with inputs permuted per fit. The DFG columns are the Part 6 runs, with inputs in natural order except on \code{dlheld2}. The last column is $\log_2$ of the portfolio's $\nstar$ over \code{nets_best}'s, which equals \code{nets_wd} on these four suites.

\begin{table}[h]
\centering\footnotesize\setlength{\tabcolsep}{4pt}
\begin{tabular}{llrrrrrr}
\toprule
target & definition & DFG v1 & portfolio & \texttt{lin} & \texttt{nets} & \texttt{nets\_wd} & $\log_2$\\
\midrule
and3 & AND$(x_1,x_2,x_3)$ & 101 & 25 & 68 & 94 & 70 & $-1.51$\\
andxor & AND$(x_1x_2,x_3x_4)$ & 43 & 40 & never & 376 & 204 & $-2.34$\\
dict & $x_1$ & 22 & 16 & 27 & 47 & 40 & $-1.33$\\
dlist5 & decision list on 5 literals & 61 & 32 & 85 & 110 & 85 & $-1.41$\\
dnf3 & 3-term DNF, 3 literals per term & 220 & 147 & never & 505 & 353 & $-1.27$\\
maj3 & MAJ$(x_1,x_2,x_3)$ & 57 & 31 & 56 & 112 & 91 & $-1.53$\\
maj9 & MAJ$(x_1..x_9)$ & 107 & 88 & 80 & 263 & 149 & $-0.75$\\
majpar & MAJ$(x_1x_2,x_3x_4,x_5x_6)$ & 62 & 60 & never & 1380 & 1026 & $-4.08$\\
mux & $x_4\,?\,x_1x_2x_3:\mathrm{MAJ}(x_1,x_2,x_3)$ & 65 & 77 & never & 1156 & 729 & $-3.25$\\
par2 & $x_1x_2$ & 16 & 16 & never & 200 & 109 & $-2.77$\\
par3 & $x_1x_2x_3$ & 16 & 55 & never & 843 & 443 & $-3.00$\\
par4 & $x_1x_2x_3x_4$ & 21 & 56 & never & 2728 & 1726 & $-4.95$\\
parmaj & MAJ$_3\oplus$MAJ$_3$ on $x_{1..6}$ & 221 & 214 & never & 743 & 434 & $-1.02$\\
rand0 & random 6-gate formula & 23 & 16 & never & 222 & 125 & $-2.97$\\
rand1 & random 6-gate formula & 104 & 94 & never & 463 & 316 & $-1.75$\\
rand2 & random 6-gate formula & 28 & 16 & 50 & 91 & 59 & $-1.89$\\
rand3 & random 6-gate formula & 189 & 181 & never & 1016 & 692 & $-1.93$\\
rmaj9 & MAJ of three MAJ$_3$ & 167 & 163 & never & 2027 & 1273 & $-2.97$\\
sparsepoly & sign$(2x_1+x_2x_3+x_4x_5x_6)$ & 73 & 95 & never & 2205 & 1238 & $-3.71$\\
stair5 & sign$(x_1+x_1x_2+\dots+x_1\cdots x_5)$ & 112 & 109 & never & 1514 & 871 & $-2.99$\\
thr10 & sign$(w\cdot x_{1..10})$, Gaussian $w$ & 164 & 121 & 113 & 563 & 307 & $-1.34$\\
\bottomrule
\end{tabular}

\caption{Toy suite \code{designed} (development; in-sample for DFG and for the Part 8 designs).}
\label{tab:appdesigned}
\end{table}

\begin{table}[h]
\centering\footnotesize\setlength{\tabcolsep}{4pt}
\begin{tabular}{lrrrrrrr}
\toprule
target & DFG v1 & portfolio & \texttt{lin} & \texttt{hgb} & \texttt{nets} & \texttt{nets\_wd} & $\log_2$\\
\midrule
rf00 & 28 & 40 & 83 & 60 & 98 & 62 & $-0.62$\\
rf01 & 48 & 48 & never & 120 & 467 & 256 & $-2.42$\\
rf02 & 51 & 49 & never & 228 & 266 & 191 & $-1.97$\\
rf03 & 87 & 71 & never & 139 & 425 & 276 & $-1.97$\\
rf04 & 54 & 48 & never & 111 & 342 & 203 & $-2.07$\\
rf05 & 74 & 76 & never & 250 & 873 & 493 & $-2.70$\\
rf06 & 27 & 16 & 37 & 58 & 61 & 50 & $-1.64$\\
rf07 & 105 & 97 & never & 168 & 407 & 281 & $-1.53$\\
rf08 & 52 & 21 & 69 & 53 & 85 & 62 & $-1.54$\\
rf09 & 57 & 45 & never & 292 & 320 & 182 & $-2.01$\\
rf10 & 56 & 36 & 66 & 53 & 103 & 79 & $-1.15$\\
rf11 & 16 & 56 & never & 416 & 864 & 445 & $-3.00$\\
rf12 & 93 & 90 & never & 524 & 963 & 610 & $-2.76$\\
rf13 & 53 & 48 & never & 213 & 401 & 245 & $-2.33$\\
rf14 & 56 & 53 & never & 378 & 714 & 404 & $-2.92$\\
rf15 & 51 & 22 & 67 & 53 & 108 & 75 & $-1.78$\\
rf16 & 54 & 36 & never & 214 & 354 & 200 & $-2.46$\\
rf17 & 28 & 56 & never & never & 3049 & 1715 & $-4.94$\\
rf18 & 51 & 46 & never & 93 & 174 & 112 & $-1.29$\\
rf19 & 56 & 52 & never & 366 & 902 & 654 & $-3.65$\\
\bottomrule
\end{tabular}

\caption{Toy suite \code{random}: 20 random formulas (development data for DFG). Networks: 2 seeds.}
\label{tab:apprandom}
\end{table}

\begin{table}[h]
\centering\footnotesize\setlength{\tabcolsep}{3.5pt}
\begin{tabular}{llrrrrrrrr}
\toprule
target & definition & DFG v1$^-$ & DFG v1.1 & portfolio & \texttt{lin} & \texttt{hgb} & \texttt{nets} & \texttt{nets\_wd} & $\log_2$\\
\midrule
hrf0 & random formula, 6 inputs & 68 & 86 & 46 & 191 & 63 & 192 & 133 & $-1.54$\\
hrf1 & random formula, 6 inputs & 117 & 118 & 117 & never & 461 & 1203 & 789 & $-2.75$\\
hrf2 & random formula, 6 inputs & 44 & 46 & 54 & 93 & 60 & 107 & 90 & $-0.74$\\
inter2 & AND of two halfspaces on $x_{1..10}$ & 704 & 583 & 394 & never & 416 & 859 & 454 & $-0.20$\\
maj15 & MAJ$(x_1..x_{15})$ & 132 & 130 & 118 & 110 & 1278 & 355 & 210 & $-0.83$\\
majmaj5 & MAJ of three MAJ$_5$ & 1252 & 816 & 1161 & never & 1638 & 2847 & 2355 & $-1.02$\\
majpar5 & MAJ$(x_1x_2,\dots,x_9x_{10})$ & 141 & 123 & 127 & never & 4127 & 5229 & 4648 & $-5.20$\\
majxor & MAJ$_5\oplus$MAJ$_5$ on $x_{1..10}$ & 782 & 766 & 771 & never & 4076 & 1313 & 792 & $-0.04$\\
mux6 & 6-multiplexer & 197 & 197 & 123 & never & 188 & 481 & 358 & $-1.54$\\
par5 & $x_1\cdots x_5$ & never & 56 & 839 & never & never & 7245 & 7096 & $-3.08$\\
stair6 & sign$(x_1+\dots+x_1\cdots x_6+\tfrac12)$ & 212 & 207 & 180 & never & 695 & 1967 & 1495 & $-3.06$\\
thr30 & sign$(w\cdot x)$, Gaussian $w$ on all 30 & 191 & 158 & 149 & 148 & 1645 & 528 & 311 & $-1.06$\\
\bottomrule
\end{tabular}

\caption{Toy suite \code{dlheld} (held-out 1). DFG v1$^-$ is frozen v1 without its degree-5/6 search, the blind run; DFG v1.1 is post hoc.}
\label{tab:appdlheld}
\end{table}

\begin{table}[h]
\centering\footnotesize\setlength{\tabcolsep}{4pt}
\begin{tabular}{llrrrrrrr}
\toprule
target & definition & DFG v1.1 & portfolio & \texttt{lin} & \texttt{hgb} & \texttt{nets} & \texttt{nets\_wd} & $\log_2$\\
\midrule
dnf4x2 & $(x_1\wedge x_2)\vee\dots\vee(x_7\wedge x_8)$ & 235 & 112 & never & 98 & 440 & 306 & $-1.45$\\
h2rf0 & random formula, 7 inputs & 117 & 108 & never & 220 & 527 & 389 & $-1.85$\\
h2rf1 & random formula, 7 inputs & 85 & 31 & 103 & 58 & 127 & 103 & $-1.73$\\
maj21 & MAJ$(x_1..x_{21})$ & 134 & 122 & 116 & 1956 & 413 & 216 & $-0.83$\\
sym8 & $[|x_1+\dots+x_8|\ge4]$ & 395 & 401 & never & 2689 & 504 & 377 & $+0.09$\\
thr20 & sign$(w\cdot x_{1..20})$, Gaussian $w$ & 136 & 130 & 117 & 1262 & 593 & 304 & $-1.23$\\
xmaj & $x_1x_2\,$MAJ$(x_3,x_4,x_5)$ & 61 & 63 & never & 1538 & 1552 & 871 & $-3.80$\\
\bottomrule
\end{tabular}

\caption{Toy suite \code{dlheld2} (held-out 2; blind for DFG v1.1, inputs permuted).}
\label{tab:appdlheld2}
\end{table}

\clearpage
\begin{table}[p]
\centering\footnotesize\setlength{\tabcolsep}{5pt}
\begin{tabular}{lrrrrrrrr}
\toprule
dataset & $n$ & MLP & MLP-L2 & \texttt{dfg\_bits} & \texttt{dfg\_gates} & \texttt{dfg\_dirs} & linear & \texttt{hgb}\\
\midrule
mnist & 300 & 0.186 & -- & 0.203 & 0.177 & 0.177 & 0.186 & 0.198\\
 & 3000 & 0.075 & -- & 0.101 & 0.095 & 0.098 & 0.113 & 0.063\\
 & 60000 & 0.025 & -- & 0.072 & 0.071 & 0.055 & 0.078 & 0.022\\
\addlinespace[1pt]
fashion\_mnist & 300 & 0.252 & -- & 0.241 & 0.318 & 0.249 & 0.257 & 0.272\\
 & 3000 & 0.169 & -- & 0.181 & 0.184 & 0.182 & 0.210 & 0.154\\
 & 60000 & 0.110 & -- & 0.142 & 0.151 & 0.156 & 0.165 & 0.105\\
\addlinespace[1pt]
adult & 300 & 0.170 & 0.181 & 0.177 & 0.177 & 0.174 & 0.168 & 0.180\\
 & 3000 & 0.155 & 0.156 & 0.153 & 0.156 & 0.152 & 0.151 & 0.147\\
 & 38842 & 0.143 & 0.140 & 0.141 & 0.141 & 0.141 & 0.146 & 0.126\\
\addlinespace[1pt]
letter & 300 & 0.379 & 0.306 & 0.375 & 0.355 & 0.356 & 0.380 & 0.348\\
 & 3000 & 0.098 & 0.094 & 0.161 & 0.138 & 0.148 & 0.235 & 0.104\\
 & 15000 & 0.033 & 0.036 & 0.069 & 0.109 & 0.114 & 0.223 & 0.041\\
\addlinespace[1pt]
pendigits & 300 & 0.054 & 0.041 & 0.092 & 0.079 & 0.075 & 0.107 & 0.100\\
 & 3000 & 0.010 & 0.010 & 0.023 & 0.018 & 0.014 & 0.059 & 0.018\\
 & 8244 & 0.006 & 0.005 & 0.013 & 0.012 & 0.010 & 0.052 & 0.011\\
\addlinespace[1pt]
connect\_4 & 300 & 0.305 & 0.306 & 0.294 & 0.300 & 0.286 & 0.294 & 0.310\\
 & 3000 & 0.229 & 0.210 & 0.243 & 0.223 & 0.241 & 0.245 & 0.215\\
 & 57557 & 0.143 & 0.141 & 0.197 & 0.179 & 0.189 & 0.236 & 0.158\\
\addlinespace[1pt]
dna & 300 & 0.144 & 0.119 & 0.060 & 0.093 & 0.101 & 0.104 & 0.057\\
 & 2390 & 0.058 & 0.048 & 0.046 & 0.052 & 0.049 & 0.058 & 0.034\\
\addlinespace[1pt]
nursery & 300 & 0.093 & 0.097 & 0.090 & 0.105 & 0.102 & 0.104 & 0.079\\
 & 3000 & 0.006 & 0.002 & 0.022 & 0.009 & 0.025 & 0.077 & 0.005\\
 & 9718 & 0.000 & 0.000 & 0.003 & 0.000 & 0.001 & 0.076 & 0.000\\
\addlinespace[1pt]
ring & 300 & 0.254 & 0.097 & 0.078 & 0.128 & 0.150 & 0.273 & 0.082\\
 & 3000 & 0.048 & 0.029 & 0.042 & 0.052 & 0.032 & 0.249 & 0.031\\
 & 5550 & 0.033 & 0.027 & 0.037 & 0.039 & 0.029 & 0.251 & 0.025\\
\addlinespace[1pt]
twonorm & 300 & 0.034 & 0.033 & 0.040 & 0.031 & 0.035 & 0.030 & 0.062\\
 & 3000 & 0.028 & 0.030 & 0.030 & 0.026 & 0.032 & 0.024 & 0.037\\
 & 5550 & 0.027 & 0.027 & 0.026 & 0.026 & 0.032 & 0.024 & 0.033\\
\addlinespace[1pt]
phoneme & 300 & 0.192 & 0.184 & 0.209 & 0.206 & 0.235 & 0.261 & 0.176\\
 & 3000 & 0.137 & 0.132 & 0.158 & 0.146 & 0.147 & 0.243 & 0.115\\
 & 4053 & 0.128 & 0.133 & 0.154 & 0.141 & 0.142 & 0.243 & 0.110\\
\addlinespace[1pt]
waveform\_40 & 300 & 0.196 & 0.180 & 0.173 & 0.157 & 0.190 & 0.181 & 0.186\\
 & 3000 & 0.149 & 0.136 & 0.136 & 0.127 & 0.148 & 0.122 & 0.141\\
 & 3750 & 0.143 & 0.127 & 0.118 & 0.120 & 0.146 & 0.119 & 0.133\\
\addlinespace[1pt]
mushroom & 300 & 0.013 & 0.012 & 0.010 & 0.006 & 0.007 & 0.022 & 0.006\\
 & 3000 & 0.002 & 0.000 & 0.000 & 0.001 & 0.001 & 0.000 & 0.000\\
 & 6093 & 0.001 & 0.000 & 0.000 & 0.000 & 0.000 & 0.000 & 0.000\\
\addlinespace[1pt]
\bottomrule
\end{tabular}

\caption{Part 7: seed-mean test error on the 13 held-out datasets at $n=300$, 3000 and the full pool. MLP-L2 was run on tabular data at $n\le3000$. On images at $n\ge30000$ the DFG adaptations are seed 0 only. \code{hgb} on Fashion-MNIST was run in Part 8. Source: \repo{study/standard_tasks/results/}.}
\label{tab:apppart7}
\end{table}

\begin{table}[p]
\centering\footnotesize\setlength{\tabcolsep}{4pt}
\begin{tabular}{lrrrrrr}
\toprule
dataset & cells & vs \texttt{mlp\_best} (b/w) & vs \texttt{mlp\_strong} (b/w) & \texttt{hgb} vs \texttt{mlp\_best} & full $n$ & error: \texttt{mlp\_best} / portfolio\\
\midrule
\multicolumn{7}{l}{\emph{HELDOUT2 (fresh)}}\\
texture & 5 & $+0.61$ (0/0) & $+0.36$ (0/0) & $+1.80$ & 4125 & 0.0027 (mlp\_strong) / 0.0061\\
churn & 5 & $-0.35$ (2/0) & $-0.36$ (2/0) & $-0.11$ & 3750 & 0.0469 (mlp\_strong) / 0.0363\\
wine\_quality\_white & 5 & $-0.15$ (2/0) & $-0.19$ (2/0) & $-0.10$ & 3674 & 0.3881 (mlp\_l2) / 0.3216\\
page\_blocks & 5 & $-0.02$ (0/0) & $-0.03$ (0/0) & $+0.40$ & 4105 & 0.0222 (mlp\_strong) / 0.0222\\
clean2 & 5 & $-1.12$ (1/0) & $-1.92$ (3/0) & $-1.14$ & 4948 & 0.0000 (mlp\_l2) / 0.0000\\
led24 & 4 & $-0.14$ (1/0) & $-0.15$ (1/0) & $-0.00$ & 2400 & 0.2925 (mlp\_l2) / 0.2933\\
ann\_thyroid & 5 & $-0.77$ (4/0) & $-0.78$ (4/0) & $-0.68$ & 5400 & 0.0176 (mlp\_l2) / 0.0074\\
banana & 5 & $-0.00$ (0/0) & $-0.04$ (0/0) & $+0.26$ & 3975 & 0.0958 (mlp\_strong1) / 0.0974\\
krkopt & 6 & $+0.19$ (0/2) & $+0.19$ (0/2) & $+0.31$ & 21042 & 0.1333 (mlp\_strong) / 0.2037\\
sleep & 8 & $-0.05$ (0/0) & $-0.10$ (3/0) & $+0.04$ & 95908 & 0.2301 (mlp\_strong) / 0.2243\\
\addlinespace[2pt]
\multicolumn{7}{l}{\emph{Part 7 tabular set (references known to the designers)}}\\
adult & 7 & $-0.11$ (2/0) & $-0.20$ (4/0)$^\dagger$ & $-0.05$ & 38842 & 0.1401 (mlp\_l2) / 0.1280\\
letter & 6 & $-0.13$ (1/0) & $-0.15$ (2/0)$^\dagger$ & $+0.21$ & 15000 & 0.0327 (mlp) / 0.0259\\
pendigits & 5 & $+0.19$ (0/0) & $-1.28$ (5/0) & $+0.99$ & 8244 & 0.0053 (mlp\_l2) / 0.0072\\
connect\_4 & 7 & $+0.03$ (0/1) & $-0.00$ (0/0)$^\dagger$ & $+0.10$ & 57557 & 0.1412 (mlp\_l2) / 0.1566\\
dna & 4 & $-0.77$ (3/0) & $-0.81$ (3/0) & $-0.60$ & 2390 & 0.0402 (mlp\_strong) / 0.0281\\
nursery & 5 & $+0.70$ (0/2) & $+0.64$ (0/2) & $+0.34$ & 9718 & 0.0000 (mlp\_l2) / 0.0007\\
ring & 5 & $-0.50$ (2/0) & $-0.94$ (3/0) & $-0.05$ & 5550 & 0.0268 (mlp\_l2) / 0.0254\\
twonorm & 5 & $-0.08$ (0/0) & $-0.15$ (0/0) & $+0.73$ & 5550 & 0.0270 (mlp) / 0.0254\\
phoneme & 5 & $-0.24$ (2/0) & $-0.27$ (2/0) & $-0.09$ & 4053 & 0.1256 (mlp\_strong) / 0.0834\\
waveform\_40 & 5 & $-0.21$ (2/0) & $-0.31$ (3/0) & $+0.09$ & 3750 & 0.1269 (mlp\_l2) / 0.1149\\
mushroom & 5 & $-0.31$ (0/0) & $-0.36$ (0/0) & $-0.14$ & 6093 & 0.0000 (mlp\_l2) / 0.0000\\
\addlinespace[2pt]
\multicolumn{7}{l}{\emph{images (references known to the designers)}}\\
mnist & 7 & $+0.08$ (1/2) & $+0.08$ (1/2) & $+0.40$ & 60000 & 0.0143 (mlp\_strong) / 0.0198\\
fashion\_mnist & 7 & $+0.03$ (0/2) & $+0.03$ (0/2) & $+0.11$ & 60000 & 0.0966 (mlp\_strong) / 0.1072\\
\addlinespace[2pt]
\bottomrule
\end{tabular}

\caption{Part 8: \code{hyb_portfolio} per held-out real dataset. ``vs \code{mlp_best}'' is the mean over cells of $\log_2(\text{error}/\text{reference error})$, followed by (better/worse) resolved cells; likewise ``vs \code{mlp_strong}''. $^\dagger$\code{mlp_strong} covers only the cells with $n\le10000$. The last four columns give the full training pool, the seed-mean error there of \code{mlp_best}, the variant that attains it (without the \code{mlp_} prefix), and the portfolio's error.}
\label{tab:apppart8}
\end{table}
\clearpage

\paragraph{Part 9 held-out toy suites.} Columns as in \cref{tab:p9dlheld6}: v2 is \code{hyb_portfolio2} and v1 \code{hyb_portfolio}, both with inputs permuted per fit; \code{nets_best} is the per-target minimum over the network variants run on the suite. These two suites are exposed to the sum-feature design (\cref{sec:p9audits}) and are secondary evidence. Definitions: \cref{sec:p9data} and \repo{study/part9/PART9.md}, \S2.2.

\begin{table}[h]
\centering\scriptsize\setlength{\tabcolsep}{3.5pt}
\begin{tabular}{lrrrrrrrrrr}
\toprule
target & v2 & v1 & \texttt{lin} & \texttt{hgb} & \texttt{nets} & \hd{\texttt{nets}\\\texttt{\_wd}} & \hd{\texttt{nets}\\\texttt{\_cewd}} & \hd{\texttt{nets}\\\texttt{\_best}} & \hd{$\log_2$ v2 /\\\texttt{nets\_best}} & \hd{$\log_2$\\v2 / v1}\\
\midrule
\texttt{h5s0} & 205 & 11042 & never & never & 1454 & 803 & 867 & 803 & $-1.97$ & $-5.75$\\
\texttt{h5s1} & never & never & never & never & never & never & never & never & both never & both never\\
\texttt{h5s2} & 1154 & 4409 & never & 10860 & 3322 & 1865 & 1850 & 1850 & $-0.68$ & $-1.93$\\
\texttt{h5s3} & 1757 & never & never & never & 8160 & 6779 & 6454 & 6454 & $-1.88$ & $-\infty$\\
\texttt{h5s4} & 184 & 415 & never & 1465 & 2744 & 1801 & 1891 & 1801 & $-3.29$ & $-1.18$\\
\texttt{h5s5} & 421 & never & never & never & 13832 & 13610 & 12254 & 12254 & $-4.86$ & $-\infty$\\
\texttt{h5s6} & 573 & never & never & never & 6227 & 4771 & 3963 & 3963 & $-2.79$ & $-\infty$\\
\texttt{h5s7} & 6760 & never & never & never & 13863 & 12878 & 7641 & 7641 & $-0.18$ & $-\infty$\\
\texttt{h5t0} & 104 & 146 & never & 2886 & 770 & 459 & 352 & 352 & $-1.75$ & $-0.48$\\
\texttt{h5t1} & 95 & 116 & never & 993 & 683 & 393 & 325 & 325 & $-1.78$ & $-0.30$\\
\midrule
\texttt{h5rf0} (control) & 268 & 311 & never & 2658 & 2673 & 2056 & 1547 & 1547 & $-2.53$ & $-0.22$\\
\texttt{h5rf1} (control) & 98 & 101 & never & 9281 & 3352 & 2859 & 2512 & 2512 & $-4.69$ & $-0.05$\\
\bottomrule
\end{tabular}

\caption{Toy suite \code{dlheld5} (held out; drawn after the sum-feature design froze, from a family inside its representable class).}
\label{tab:appdlheld5}
\end{table}

\begin{table}[h]
\centering\scriptsize\setlength{\tabcolsep}{3.5pt}
\begin{tabular}{lrrrrrrrrrr}
\toprule
target & v2 & v1 & \texttt{lin} & \texttt{hgb} & \texttt{nets} & \hd{\texttt{nets}\\\texttt{\_wd}} & \hd{\texttt{nets}\\\texttt{\_cewd}} & \hd{\texttt{nets}\\\texttt{\_best}} & \hd{$\log_2$ v2 /\\\texttt{nets\_best}} & \hd{$\log_2$\\v2 / v1}\\
\midrule
\texttt{band10} & 432 & 3468 & never & never & 5920 & 4722 & 3440 & 3440 & $-2.99$ & $-3.00$\\
\texttt{ex4of9} & 194 & 1658 & never & never & 1073 & 741 & 411 & 411 & $-1.08$ & $-3.09$\\
\texttt{majsym} & 189 & 352 & never & 2840 & 1542 & 1062 & 595 & 595 & $-1.66$ & $-0.90$\\
\texttt{ssym8} & 157 & 1410 & never & 8072 & 1160 & 715 & 396 & 396 & $-1.34$ & $-3.17$\\
\texttt{sumpar} & 164 & 168 & never & never & 3487 & 2619 & 1837 & 1837 & $-3.48$ & $-0.04$\\
\texttt{sym11} & 203 & 2386 & never & never & 934 & 656 & 374 & 374 & $-0.88$ & $-3.56$\\
\texttt{symand} & 170 & 289 & never & 1307 & 498 & 382 & 370 & 370 & $-1.12$ & $-0.77$\\
\texttt{symxor} & 403 & 766 & never & 6909 & 1706 & 946 & 969 & 946 & $-1.23$ & $-0.93$\\
\texttt{tail13} & 176 & 436 & never & never & 612 & 458 & 420 & 420 & $-1.26$ & $-1.31$\\
\midrule
\texttt{dnf3x3} (control) & 118 & 118 & never & 148 & 493 & 373 & 392 & 373 & $-1.66$ & $+0.00$\\
\texttt{h4rf0} (control) & 48 & 48 & never & 160 & 349 & 206 & 240 & 206 & $-2.10$ & $+0.00$\\
\texttt{h4rf1} (control) & 56 & 51 & 158 & 61 & 211 & 143 & 181 & 143 & $-1.35$ & $+0.14$\\
\bottomrule
\end{tabular}

\caption{Toy suite \code{dlheld4} (held out, fixed before any Part 9 design; its definitions were seen by the sum-feature builder).}
\label{tab:appdlheld4}
\end{table}
